\begin{textsample}{NEG Dim 9 – global\_south\_2024\_11 – Score -3.00 – global\_south\_2024\_11/117424852.txt}  \label{ex:f9_neg_419}
Guests of Terror : Geopolitics and the Apolitical Literati of the Philippines A handful of Filipino authors made public their decision to refuse the horror"honor"of being invited to the Frankfurt Book Fair ( FBF ), especially after said event declared its Zionist position of standing"with complete solidarity on the side of Israel"back in October 2023, punctuated by the indefinite delay and eventual cancellation of the awarding ceremony for Palestinian author Adania Shibli, and her novel Minor Detail.
This"ethically indefensible decision"in effect censored a literary work that"juxtaposes the true story of the rape and \textbf{murder} of a Bedouin \textbf{girl} by an Israel army unit in 1949 with the fictional story of a \textbf{female} journalist investigating the crime in the Palestinian city of Ramallah."Not content with a low ( and somewhat defensive ) blow against freedom of expression, FBF director Juergen Boos launched a predictable scripted offensive and decried the crocodile-tear-inducing accusatory"anti-Semitism."He also announced the staging in the Fair ' s cultural and @ @ @ @ @ @ @ @ @ @"Out of Concern for Israel"The genocide-complicit Frankfurt Book Fair shrugs off what UN experts called a"scholasticide"that left an unprecedented trail of blood and rubble, of academics and research infrastructure.
Israel targeted and killed Sufian Tayeh, leading researcher in theoretical physics and applied mathematics and the president of the Islamic University of Gaza, along with some of his family members.
Sufian Tayeh is just one of the many Palestinian scientists, academic scholars, and artists killed by Israel as part of the occupation entity ' s scholasticide campaign.
But Filipino writers attending the FBF -- some of whom filled slots vacated by those who declined the Zionist-wrath-farmed laurel wreath -- turn a blind eye and/or cope by convincing themselves and their projected national readership that visibility in a renowned international platform is a much-needed break to introduce Philippine literature to the world.
Eyeing the"guest of honor"status of the country come 2025, the apolitical literati plays the dangerous game of superficial representation politics as means to promote homegrown"genius @ @ @ @ @ @ @ @ @ @ a dispossessed people ( such as ours ) relegates another dispossessed people ( such as the Palestinians ) to collateral irrelevance, in contradistinction to the aspirational call"From Palestine to Philippines, Stop the US War Machine"-- a call of solidarity that might be unfamiliar or worse, deemed"anti-Semitic"and"violent"by FBF enjoyers, echoing their un/intended Zionist patrons.
Boasting of more than 60,000 new book titles in 2023 alone, Germany rightly deserves the reputation of hosting ( and regulating ) the world ' s most celebrated book fair.
But as the **32 ; 2373 ; TOOLONG clich? goes, FBF is again far from being fair, serving as bad breath of air freshener overpowered by the tang of burnt charcoal and saltpeter that rises from imperialist-prompted war zones in Syria, Lebanon, Iran, Iraq, Palestine, and Lebanon.
The Frankfurt Book Fair, as an institution has never once distanced itself from, nor is it critical of, Germany ' s avid role in and unconditional support of the Palestinian genocide.
On @ @ @ @ @ @ @ @ @ @, expressing uncritical support for Israel, said that"Israel can kill civilians".
And the statement was made around the same time period as the release of a video of a Palestinian hospital patient, a nineteen year-old teen who was shown still connected to an IV drip and flailing helplessly while being burned to death by an Israeli airstrike.
Not being a Palanca awardee hence never on the way to Hall of Fame breakthrough then National Artist rockstardom then towards the wettest of the wet dreams of Philippine literati : the Nobel Peace prize world-domination, I am ( objectively speaking, esp. in stark contrast to the two aforementioned greats ) a nobody who is just somebody preoccupied ( even fixated ) with matters other than literary? clat.
This article is driven by a simple and basic ( not really -- or rather far from --"radical") humanist position that the likes of Dalisay and Almario might ( feign to? ) share with the rest of the nation and the world.
Isn't the"@ @ @ @ @ @ @ @ @ @ criticism?
Critics of the Palanca, who likewise deserve to be heard, remain voices in the wilderness that reach a few people compared with the reach of award-winners and wet-dreamers.
Isn't the supposed"nationalist"stance to be in solidarity with nations in the same boat as ours?
What was thought as fence-sitting and both-sides-ing is tilting the power dynamics in favor of the dominant positionality -- a"centrism"that is rather illiberal and faux-nationalist.
Dogmatic pacifism amid genocide is closet fascism. *** Resonant with Junior ' s Marcosian Bagong Pilipinas, the surface-level apolitical participation of Filipino writers in a genocide-complicit literary affair such as the FBF champions a hegemonic politics that privileges careeristic survival over basic humanist compassion.
Does grabbing the opportunity to promote novels, anthologies of short stories, collections of essays, chapbooks of poetry, komiks, picture books for children, take precedence over censorship of author by virtue of race, lives of fellow writers cut short hence denying them of writing more books, children robbed of their @ @ @ @ @ @ @ @ @ @ libraries and schools?
What to do with literary careers that licked the blood off of the hands of genocidal maniacs worse than perpetrators of drug war and cronies of martial law?
Isn't it time to bite?
No?
Are congratulations still in order as courtesy to our esteemed writers, envied by naysayers, greatest wordsmiths of the Malayan race, the best offerings of Philippine letters?
Mabuhay?
Padayon?
Davao Today is a daily online news magazine based in Davao city.
News events concerning business and politics in the city and nearby towns and cities are covered by Davao Today news team on a daily basis.
Davao Today offers features on lifestyle and culture, consumers, agriculture, education, and the lives of hardworking Davaoe? os.
Copyright 2017 DAVAOTODAY.com All Rights Reserved Terms of Use Advertise with Us Contact Us

% matched lemmas: female, girl, murder
\end{textsample}
